\documentclass[11pt]{amsart}
\usepackage[margin=1in]{geometry}
\usepackage{amsmath,amssymb,amsthm,mathtools}
\usepackage{graphicx}
\usepackage[hidelinks]{hyperref}
\usepackage{microtype}

\newtheorem{theorem}{Theorem}
\newtheorem{corollary}[theorem]{Corollary}
\newtheorem{lemma}[theorem]{Lemma}
\theoremstyle{remark}
\newtheorem{remark}[theorem]{Remark}

\newcommand{\R}{\mathbb R}
\newcommand{\mdim}{\operatorname{mdim}}
\newcommand{\relint}{\operatorname{relint}}

\title{Sharp metric-dimension results around Geh\'er's Problem 1}
\author{Automated proof attempt for arXiv:1507.05114}
\date{August 13, 2026}

\begin{document}
\begin{abstract}
For a $d$-dimensional normed space, we prove that every finite resolving set
has at least $d+1$ points.  Norms whose unit ball has a facet have no finite
resolving set at all.  At the opposite end, the metric dimension of
$\ell_p^d$ is exactly $d+1$ for every $1<p<\infty$, with the standard simplex
as an optimal resolving set.  We also prove the same conclusion for a
quantitative $C^{1,1}$ neighbourhood of Euclidean norms containing genuinely
non-Hilbert examples.  These results answer the literal non-strict question
negatively and settle large strictly convex classes, but do not resolve the
universal strictly convex version.
\end{abstract}
\maketitle

\section{Source question and scope}

Let $(X,\|\cdot\|)$ be a finite-dimensional real normed space.  A set
$R\subset X$ \emph{resolves} $X$ if
\[
 \bigl(\|x-r\|\bigr)_{r\in R}=\bigl(\|y-r\|\bigr)_{r\in R}
 \quad\Longrightarrow\quad x=y.
\]
We write $\mdim X$ for the least cardinality of a finite resolving set, and
$\mdim X=\infty$ if none exists.

Section 4, Problem 1 of Geh\'er~\cite{Geher} asks exactly:
\begin{quote}
``Is the metric dimension of every $d$-dimensional (strictly convex) normed
space less than or equal to $d+1$?''
\end{quote}
The next paragraph says that it was not clear whether the metric dimension
could be smaller than $d+1$.  The exact source context is reproduced below.

\begin{center}
\includegraphics[width=.96\textwidth]{figures/open_problem_crop.png}
\end{center}

The parenthetical admits two natural readings: all normed spaces, or at least
all strictly convex ones.  The first is false by Theorem~\ref{thm:facet}.  The
second is verified here for broad classes but remains unresolved in full.

\section{Proof intuition}

There are three short mechanisms.
\begin{enumerate}
\item Distances to $m$ landmarks give a continuous coordinate map into
$\R^m$.  Invariance of domain forbids injectivity when $m\le d$.
\item A facet creates an open cone on which the norm is exactly linear.  Far
inside that cone, a finite landmark set cannot detect motion tangent to the
facet.
\item In $\ell_p^d$, subtracting the equality at the origin from the equality
at the $i$th coordinate landmark isolates one strictly monotone scalar
function of the $i$th coordinate.  Near a Hilbert norm, the same relative
distance map is an invertible linear map plus a strict Lipschitz contraction.
\end{enumerate}

\section{A universal sharp lower bound}

\begin{theorem}[No $d$ landmarks]\label{thm:lower}
Let $X$ be a real normed space of dimension $d\ge1$.  If $X$ has finite metric
dimension, then
\[
 \mdim X\ge d+1.
\]
\end{theorem}

\begin{proof}
Identify $X$ linearly and topologically with $\R^d$.  Suppose that
$R=\{r_1,\dots,r_m\}$ resolves $X$, where $m\le d$, and define
\[
 F:X\longrightarrow\R^m,\qquad
 F(x)=(\|x-r_1\|,\dots,\|x-r_m\|).
\]
This map is continuous and injective.  Pad it with $d-m$ zero coordinates to
obtain a continuous injection $\widetilde F:\R^d\to\R^d$.  By Brouwer's
invariance-of-domain theorem, $\widetilde F(\R^d)$ is open.

If $m<d$, this image lies in a proper coordinate subspace and hence has empty
interior, a contradiction.  If $m=d$, then
$F(X)\subset[0,\infty)^d$.  Moreover, the first coordinate of $F(r_1)$ is
zero, so no Euclidean neighbourhood of $F(r_1)$ is contained in the
nonnegative orthant.  Thus $F(X)$ is not open, again a contradiction.
\end{proof}

This proves the lower bound conjectured in the paragraph following Problem 1,
without any convexity or smoothness assumption.

\section{Facets force infinite metric dimension}

Let $B$ be the closed unit ball.  By a \emph{facet} we mean an exposed face
of affine dimension $d-1$.

\begin{theorem}[Facet obstruction]\label{thm:facet}
Let $d\ge2$.  If the unit ball of $X$ has a facet, then
$\mdim X=\infty$.
\end{theorem}

\begin{proof}
Choose $\varphi\in X^*$ such that $\varphi\le1$ on $B$ and
\[
 F=B\cap\{\varphi=1\}
\]
is a facet.  The cone
\[
 C=\{tz:t>0,\ z\in\relint F\}
\]
is a nonempty open subset of $X$, and
\begin{equation}\label{eq:linearcone}
 \|z\|=\varphi(z)\qquad(z\in C).
\end{equation}

Let $R\subset X$ be finite.  Fix $q\in C$ and a nonzero
$w\in\ker\varphi$.  Since $C$ is open and conic, for all sufficiently large
$T$ and all $r\in R$ both
\[
 Tq-r\quad\hbox{and}\quad Tq+w-r
\]
belong to $C$.  Put $x=Tq$ and $y=Tq+w$.  Then $x\ne y$, while
\eqref{eq:linearcone} gives, for every $r\in R$,
\[
 \|x-r\|=\varphi(Tq-r)
 =\varphi(Tq+w-r)=\|y-r\|.
\]
Thus $R$ does not resolve $X$.  Since $R$ was arbitrary, no finite resolving
set exists.
\end{proof}

\begin{corollary}\label{cor:facet}
Every polyhedral normed space of dimension at least two has infinite metric
dimension.  Every non-strictly-convex normed plane has infinite metric
dimension.
\end{corollary}

\begin{proof}
A full-dimensional polytope has a facet.  In dimension two, every
nondegenerate segment in the unit circle extends to a one-dimensional exposed
face, which is a facet.
\end{proof}

In particular, $\ell_1^d$ and $\ell_\infty^d$ have infinite metric dimension.
This also shows why strict convexity is indispensable in the planar version.

\section{The exact value for \texorpdfstring{$\ell_p^d$}{ell-p spaces}}

\begin{theorem}[Coordinate-power spaces]\label{thm:lp}
For $d\ge1$ and $1<p<\infty$,
\[
 \mdim(\ell_p^d)=d+1.
\]
More precisely, $\{0,e_1,\dots,e_d\}$ resolves $\ell_p^d$.
The same conclusion holds for every weighted norm
\[
 \|x\|=\left(\sum_{i=1}^d c_i|x_i|^p\right)^{1/p},\qquad c_i>0,
\]
and every invertible linear image of such a norm.
\end{theorem}

\begin{proof}
It suffices to treat the weighted form.  Suppose $x,y\in\R^d$ have the same
distances to $0,e_1,\dots,e_d$.  Equality at the origin says
\begin{equation}\label{eq:base}
 \sum_j c_j|x_j|^p=\sum_j c_j|y_j|^p.
\end{equation}
Equality at $e_i$, after taking $p$th powers and subtracting
\eqref{eq:base}, gives
\[
 |x_i-1|^p-|x_i|^p=|y_i-1|^p-|y_i|^p.
\]
Set $h(t)=|t-1|^p-|t|^p$ and
$j_p(t)=\operatorname{sgn}(t)|t|^{p-1}$.  Then
\[
 h'(t)=p\bigl(j_p(t-1)-j_p(t)\bigr)<0,
\]
because $j_p$ is strictly increasing and $t-1<t$.  Hence $h$ is strictly
decreasing, so $x_i=y_i$.  This holds for every $i$, and therefore $x=y$.
Theorem~\ref{thm:lower} proves optimality.  An invertible linear isometry
transports resolving sets and metric dimension.
\end{proof}

\section{A smooth non-Hilbert neighbourhood}

The next result is useful because it does not rely on coordinate-power
separability.  All Euclidean norms and sufficiently small $C^{1,1}$
perturbations satisfy it.

\begin{theorem}[Quantitative near-Hilbert class]\label{thm:near}
Let $N$ be a norm on $\R^d$ and write
\[
 \Phi(x)=N(x)^2=|x|_2^2+H(x).
\]
Assume that $H$ is continuously differentiable and that
\[
 |\nabla H(u)-\nabla H(v)|_2\le\eta|u-v|_2
 \quad(u,v\in\R^d)
\]
for some $\eta<2/\sqrt d$.  Then $\{0,e_1,\dots,e_d\}$ resolves
$(\R^d,N)$, and consequently
\[
 \mdim(\R^d,N)=d+1.
\]
\end{theorem}

\begin{proof}
Define the relative squared-distance map $G:\R^d\to\R^d$ by
\[
 G_i(x)=\Phi(x-e_i)-\Phi(x)=1-2x_i+E_i(x),
 \qquad E_i(x)=H(x-e_i)-H(x).
\]
For $x,y\in\R^d$, the fundamental theorem of calculus on the segment from
$y$ to $x$ yields
\begin{align*}
 E_i(x)-E_i(y)
  =\int_0^1\!\langle
  &\nabla H(y+t(x-y)-e_i)\\
  &-\nabla H(y+t(x-y)),x-y\rangle\,dt.
\end{align*}
The two gradient arguments in the integrand are one unit apart.  Therefore
\[
 |E_i(x)-E_i(y)|\le\eta|x-y|_2,
 \qquad |E(x)-E(y)|_2\le\sqrt d\,\eta|x-y|_2.
\]
If $G(x)=G(y)$, then
\[
 2|x-y|_2=|E(x)-E(y)|_2
 \le\sqrt d\,\eta|x-y|_2.
\]
Since $\sqrt d\eta<2$, this forces $x=y$.  Equal distances to the origin and
the $e_i$ imply equality of all the $G_i$, so the landmarks resolve.  Apply
Theorem~\ref{thm:lower} for optimality.
\end{proof}

\begin{corollary}[Explicit non-Hilbert examples]\label{cor:example}
For every $d\ge2$, the norms
\[
 N_\varepsilon(x)
 =\bigl(|x|_2^2+\varepsilon\|x\|_4^2\bigr)^{1/2}
\]
have metric dimension $d+1$ for all sufficiently small $\varepsilon>0$.
They are smooth, strictly convex, and non-Hilbert.
\end{corollary}

\begin{proof}
The $\ell_2$-sum of two norms is a norm, and the Euclidean term makes it
strictly convex.  The function $x\mapsto\|x\|_4^2$ has globally Lipschitz
gradient: away from zero its Hessian is homogeneous of degree zero and bounded
on the Euclidean unit sphere, while the degree-one gradient extends
Lipschitzly through zero.  Hence Theorem~\ref{thm:near} applies when
$\varepsilon$ is small enough.  Finally, with $x=e_1$ and $y=e_2$, the two
sides of the parallelogram identity contain respectively the unequal
$\ell_4^2$ contributions $2\varepsilon\sqrt2$ and $4\varepsilon$.
\end{proof}

\section{Why the full strictly convex question is not claimed}

Eight focused routes were checked.  The successful ones are the lower bound,
facet cone, coordinate-power separation, and quantitative near-Hilbert
argument above.  Four further routes identify genuine obstructions:
\begin{itemize}
\item Strictly convexifying a facet norm does not preserve its collisions:
$\ell_q^d\to\ell_\infty^d$ as $q\to\infty$, yet every finite $q>1$ has metric
dimension $d+1$.
\item The topological obstruction ends exactly at $d$ target coordinates; an
extra distance coordinate removes the invariance-of-domain contradiction.
\item Strict planar bisectors are direction-dependent graphs, but their
strips are not uniform over the unknown pair.  Incidence counting does not
force a simplex avoiding every bisector.
\item At infinity, smooth strictly convex norms encode direction through the
dual Gauss map.  Without a quantitative curvature bound, uniform asymptotic
approximation does not prevent folds at finite or intermediate scale.
\end{itemize}
Thus the packet neither proves nor disproves
$\mdim X=d+1$ for every strictly convex $d$-dimensional normed space.

\section{Verification, literature caution, and review}

All arguments are exact and analytic; no numerical evidence is used.  The
source was rebuilt from the cached arXiv file, and the displayed crop is taken
from source PDF page 12.  Cheap run indexes and targeted searches found no
later primary-source answer, but terminology is unusually ambiguous and older
work often says \emph{determining set} or \emph{equidistant set}.  In
particular, Kalisch--Straus~\cite{KalischStraus} should be checked by a human
specialist before any novelty claim.  No priority claim is made.

Human review should focus on the intended parenthetical scope of Problem 1,
the openness of the cone over a facet, and the global $C^{1,1}$ hypothesis in
Theorem~\ref{thm:near}.  The packet should not be cited as resolving the
universal strictly convex case.

\begin{thebibliography}{9}
\bibitem{Geher}
G. P. Geh\'er,
\emph{Is it possible to determine a point lying in a simplex if we know the
distances from the vertices?}, arXiv:1507.05114.

\bibitem{KalischStraus}
G. K. Kalisch and E. G. Straus,
\emph{On the determination of points in a Banach space by their distances from
the points of a given set}, An. Acad. Brasil. Ci. \textbf{29} (1957), 501--519.
\end{thebibliography}

\end{document}
